\begin{claim}{C.79}{Proved}{Every relative disk loop for a one-dimensional proper operator action fills}
Let $O=\operatorname{int}\overline D$ be the interior
of the actual closed disk. Let
$Q\leq\operatorname{Homeo}(O)$ be closed, locally
compact, free and proper, with $\dim Q=1$.
Use the actual relative centralizer
\[
 \mathscr C_Q=
 \{h\in\operatorname{Homeo}(\overline D):
       h|_{\partial\overline D}=\id,
       \ hq=qh\text{ on }O\text{ for every }q\in Q\}
                                                               \tag{C.79.1}
\]
with its compact-open topology on maps and inverses.
Every continuous based loop $f_\tau$ in this group
has a continuous based contraction in the same group.
In particular
\[
                         \pi_1(\mathscr C_Q,1)=0.
                                                               \tag{C.79.2}
\]
The construction uses the original disk boundary
and every original operator. Its map and inverse
families are jointly continuous there, including
arbitrary accumulation of principal strips.
It does not require a specified Fuchsian disk
comparison, and does not replace the original
compactification by another one.

The two-endpoint step retains one common increasing
base map and one compact global source family in
both endpoint selections. Their fiber displacements
are blended in that same original real fiber.
This gives an actual filling on the whole disk;
no extension through an omitted strip is left in
the stated dimension-one disk assertion.
\end{claim}

\begin{proof}
\emph{The exact strips and the common preliminary family.}
Use the actual principal coordinate and inverse
of C.75. Write
\[
        f_\tau(s,t)=(s+u_\tau(t),\beta_\tau(t)),
        \qquad A=\operatorname{Fix}\langle\beta_\tau\rangle.
                                                               \tag{C.79.3}
\]
Retain the original group $G=\langle f_\tau\rangle$
and its closure $L$ in the actual relative group
(C.79.1). The complete compatible metric, all
component operators, and the jointly continuous
principal-coordinate functions on compact source
families are exactly those in C.76--C.78.
There are countably many complementary intervals.

C.78 fills all their types except bounded intervals
$I=(a,b)$ with nonabelian effective base image and
actual joint character image
\[
                  c_I=(c_a,c_b):G\longrightarrow\mathbb R^2,
                  \qquad c_I(G)=\mathbb R^2.
                                                               \tag{C.79.4}
\]
We construct all strip fillings together, using
one compact family, so that their accumulation
estimates remain uniform. Begin with precisely
the preliminary compact $K$ and $K_n$ of C.78:
$K_0=P^5$, where $P$ is the original loop and its
inverse, followed by fifteenfold short-increment
families whose infinite products converge in $L$.
This construction is independent of the strip.

\emph{A full-vector average gives the one common base.}
On a strip (C.79.4), replace (C.78.4) by the
full two-character fibers
\[
       F_n(r,\tau)=
          \{k\in K_n:c_I(k)=r c_I(f_\tau)\},
       \qquad
       F_\infty(r,\tau)=
          \{k\in K:c_I(k)=r c_I(f_\tau)\}.
                                                               \tag{C.79.5}
\]
The planar additive identity (C.37.3) gives
nonempty fibers already in $P^3\subset P^5$.
For nearby parameters the required correction
is the actual vector
$r'c_I(f_{\tau'})-r c_I(f_\tau)$.
C.76's six-short-increment neighborhood covers
this vector when it is sufficiently small; it
is a subfamily of the fifteenfold increments
used here. The same increment corrects every
label in the given fiber. Thus the Hilbert
convex-average proof (C.78.5)--(C.78.8) applies
to (C.79.5) without change: its lower-limit
argument needs only a common arbitrarily small
character correction, not a scalar target.

Choose its continuous positive average $q_p$,
$p=(r,\tau)$, supported in the precise sense
$q_p\in\mathcal P(F_\infty(p))$, with the same
four exact parameter faces. Put
\[
               B_p(t)=\Lambda_{q_p}(k\mapsto\beta_k(t)).
                                                               \tag{C.79.6}
\]
The compact positive-gap and inverse-envelope
arguments (C.78.9)--(C.78.10) make this a jointly
continuous family of increasing homeomorphisms
of $I$, with jointly continuous inverses.
It has the original base map at $r=1$ and the
identity on the other based faces. Crucially,
this is the \emph{same} $B_p$ for both endpoints.

\emph{Two actual scalar comparison planes.}
For $e=a,b$ retain the actual source action
\[
        \begin{split}
        T_e(k)(z,t)&=(z+c_e(k),\beta_k(t)),\\
        T_e(k)^{-1}(z,t)&=(z-c_e(k),\beta_k^{-1}(t)).
        \end{split}
                                                               \tag{C.79.7}
\]
In each of these planes use $z$ as the real
target, whose fibers are the proper copies
of $I$. Every map has an increasing target
translation and an increasing fiber map.
The minimum and maximum labels for the
positive average (C.79.6) belong to
$F_\infty(p)$ and hence have \emph{both}
endpoint characters prescribed. Apply the
C.62 triangle separately in the two planes,
along a path in the same compact $K$ joining
those labels. It follows that the single
compact family
\[
                         M_0=K K^{-1}K
                                                               \tag{C.79.8}
\]
realizes $(r c_a(f_\tau),B_p(t))$ under $T_a$
at $(0,t)$, and separately realizes
$(r c_b(f_\tau),B_p(t))$ under $T_b$ there.
The realizing labels need not be equal.
No realization retaining both characters
at once is claimed at this step.

Each scalar action is transitive. Indeed,
C.77 proves that the actual source derived
subgroup has transitive base image and
vanishes under every endpoint character;
each character in (C.79.4) has full real
image. Thus the precise based path-group
orbit maps of both scalar actions are open,
with the actual kernels and stabilizers
described in (C.77.4) and (C.77.12).

Add these two comparison spaces per interval
to the countable collection in C.78. Apply
its proved compact local pair-transport
construction (C.78.13) to the whole collection.
The second infinite-product construction
(C.78.14) then yields \emph{one} compact
path-connected global family $M\subset L$
and nested compact families $M_n$ starting
with (C.79.8). This includes the comparison
spaces for all the older strip types as well.
All maps and inverses are original closed-disk
maps fixing the original circle.

\emph{The two fiber selections and their exact blend.}
For $e=a,b$, define
\[
 A_n^e(p,t)=\operatorname{conv}
 \{u_k(t):k\in M_n,
       \ c_e(k)=r c_e(f_\tau),\ \beta_k(t)=B_p(t)\}.
                                                               \tag{C.79.9}
\]
They are nonempty compact intervals by
(C.79.8). The local pair transports in their
respective scalar comparison planes prove
the interlaced lower-limit property exactly
as in (C.78.16). The four relative face
values are represented by the identity
and the original $f_\tau$. The same relative
interval selection therefore supplies two
continuous displacements $U_a(p,t)$ and
$U_b(p,t)$, each in its limiting compact
interval, with those four faces exactly.
Their actual base map is the common $B_p$.

Choose the explicit cutoff
\[
 \begin{split}
 \theta_I(t)&=\max\left(0,
          \min\left(1,3\frac{t-a}{b-a}-1\right)\right),\\
 U(p,t)&=(1-\theta_I(t))U_a(p,t)+\theta_I(t)U_b(p,t).
 \end{split}
                                                               \tag{C.79.10}
\]
It is zero on a neighborhood of $a$ and
one on a neighborhood of $b$. Define
\[
 \begin{split}
 H_p(s,t)&=(s+U(p,t),B_p(t)),\\
 H_p^{-1}(s,v)&=
        (s-U(p,B_p^{-1}(v)),B_p^{-1}(v)).
 \end{split}                                                   \tag{C.79.11}
\]
These are exact inverse homeomorphisms,
jointly continuous on the open strip.
Both selections have the same four faces,
so their blend does too.

Every value of $U_a$ or $U_b$ lies between
actual $M$-label displacements at that same
base level. Their convex blend lies in the
convex hull of the union of those labels.
The minimum and maximum of the compact
set of \emph{all} $M$-labels at that level
therefore bracket $U$. Applying the original
source-plane C.62 triangle and its inverse
bracketing gives directly
\[
           H_p(x)\in M M^{-1}M(x),\qquad
           H_p^{-1}(x)\in M^{-1}M M^{-1}(x).
                                                               \tag{C.79.12}
\]
There is no new word-family bound inferred
from an unrelated coordinate interpolation:
both selected displacements and their blend
have the identical actual real base level.
The older strip constructions use the same
$M$ and have the same two bounds by C.78.

\emph{All component operators.}
The original component generator is
$q(s,t)=(\lambda s+\kappa,t+d)$, with
$d>0$ and $\lambda\ne0$, when it exists.
For the translated interval $I+d$, the
two endpoint characters are multiplied by
$\lambda$, so the fibers (C.79.5) are the
same actual subsets of $K$. Transport the
positive average by that identical label
choice. Then $B'_p(t+d)=B_p(t)+d$, exactly
as in C.78. Transport both fiber selections
by $U'_{e+d}(p,t+d)=\lambda U_e(p,t)$.
The explicit cutoff satisfies
$\theta_{I+d}(t+d)=\theta_I(t)$, and hence
the blended displacement has the same
multiplier equation. This verifies actual
commutation with $q$, including negative
$\lambda$; the original constant $\kappa$
cancels from the equation and remains in
the operator. No nonzero component power
stabilizes a bounded interval, so transport
from interval representatives is independent
of the chosen operator. The older interval
types, including $I=\mathbb R$ with its
possible component stabilizer, retain their
exact C.78 line-bundle and transport equations.

\emph{The two finite-endpoint sequence alternatives.}
On every line over $a\in A$, prescribe
$H_p(s,a)=(s+r u_\tau(a),a)$.
Consider a sequence of source points with
$t_n\to a\in A$, allowing all parameters
and all interval indices to vary. The older
strip types have the joint limit already
proved in C.78, with the same compact $M$.
For the new bounded strips, every subsequence
has a further subsequence of one of these
two forms.

First, one fixed interval $I=(a_0,b_0)$
occurs throughout. The limiting point of
$A$ must be $a_0$ or $b_0$.
Near $a_0$ the cutoff is identically zero,
so $U=U_a$ and every defining label has
$u_k(a_0)=r u_\tau(a_0)$.
Near $b_0$ it is identically one, and the
same statement uses $U_b$ and $b_0$.
Uniform continuity over the compact actual
$M$ proves the required limit of $U$.

Second, the bounded intervals
$I_n=(a_n,b_n)$ are pairwise distinct.
Then both $a_n$ and $b_n$ tend to $a$.
For example, if along a subsequence
$b_n\geq a+2\epsilon$ while $t_n\to a$,
then for all sufficiently large $n$ the
fixed point $a+\epsilon$ lies in $I_n$,
contradicting disjointness of distinct
components. The analogous argument with
$a-\epsilon$ handles the other endpoints.
Thus uniform continuity of the one compact
family $M$ gives, respectively,
\[
 \begin{split}
 U_a(p_n,t_n)-r_nu_{\tau_n}(a_n)&\longrightarrow0,\\
 U_b(p_n,t_n)-r_nu_{\tau_n}(b_n)&\longrightarrow0.
 \end{split}                                                   \tag{C.79.13}
\]
For the first estimate use the labels in
the limiting interval of $A_n^a$, and for
the second use those for $A_n^b$; they
need not be the same labels. Each estimate
persists under its convexification.
The original compact loop makes both
endpoint terms approach the prescribed
value at $a$, uniformly in the parameters.
Any convex combination in (C.79.10) has
the same limit. This proves the desired
limit even when the cutoff values vary
without converging. The two subsequence
alternatives establish the assertion for
the original arbitrary sequence.

The base average and its inverse satisfy
the compact envelopes (C.78.17), as do
the older prescriptions. Therefore their
values tend to $a$ uniformly whenever
the source or inverse input tends to $a$.
For the inverse, put $v_n=B_{p_n}^{-1}(t_n)$.
Then $v_n\to a$ and it stays in the same
interval $I_n$. Apply the same two
alternatives at $v_n$ in (C.79.11).
This gives the negative prescribed fiber
displacement and proves inverse continuity
at the limiting line. Along $A$ itself,
both limits follow from the original
continuous function $r u_\tau(a)$.

\emph{The actual whole-disk filling.}
Use (C.79.11) on every new strip, the
simultaneously constructed C.78 maps on
every other strip, the prescribed map
on each line over $A$, and the identity
on the original boundary circle. These
form a bijection of the actual closed
disk with the displayed piecewise inverse,
since every base interval and every line
over $A$ is preserved. Interior joint
continuity of maps and inverses is exactly
the strip and endpoint comparison just
proved.

On a line over $A$, the scalar fraction
$r u_\tau(a)$ is attained by an original
prefix value; its negative is attained
by the corresponding inverse. Thus the
compact bounds (C.79.12) hold there too.
Both global word families in (C.79.12)
are compact families of actual closed-disk
homeomorphisms fixing the original circle.
As a point approaches any original circle
point, all their evaluations approach that
same point uniformly. This proves joint
boundary continuity of the pasted maps
and their inverses, including arbitrarily
varying strips and unbounded principal
coordinates.

Hence the family consists of actual
homeomorphisms of $\overline D$ and is
continuous in the topology of (C.79.1).
Its operator equations were checked for
every actual component and identity-component
operator. Its four parameter faces are
$H_{0,\tau}=1$, $H_{1,\tau}=f_\tau$,
and $H_{r,0}=H_{r,1}=1$.
This is the asserted based contraction,
proving (C.79.2).

The proof is a loop statement for the
specified actual disk group. It asserts
neither contractibility of that entire
group nor a homotopy lift through C.73's
prime-end injection for a different
original domain. Those exact comparison
and broader C.12 requirements remain
as previously stated.
\end{proof}
